\section{An Elementary Theory of the Category of Sets}
\label{sec:etcs}

With a deterministic, structurally-typed substrate established in Section~2, mathematical sets can be bootstrapped inside the type system. Here, we aim to model both the syntax and the semantics of mathematical set-builder notation in executable \textbf{Jlt}, hoping to achieve a system that lends itself to the expression and subsequent evaluation of precisely those kinds of problems that one might also encounter in a pure math blackboard setting. 

A common implementation approach follows material Zermelo-Fraenkel (ZF) set theory, where the discussion starts from an element membership relation $\in$ acting on a cumulative universe of sets. Implementations such as \texttt{std::set} use this as a starting point and model sets as \textit{extensional memory containers}---meaning the identity of a set is determined by explicitly enumerating its members in physical allocations, and membership is evaluated by traversing a data structure (e.g., a balanced binary search tree). This approach exhibits prohibitive shortcomings when applied to algebraic and engineering problems: for example, it cannot natively model infinite, un-enumerable sets like the natural numbers $\mathbb{N}$, or un-enumerable continuous intervals over the real numbers $\mathbb{R}$. As a result, both the structural guarantees a compiler can infer and the optimizations it can perform on such problems are severely restricted.

By contrast, the \textit{structural categorical comprehension} view pioneered by William Lawvere~\cite{lawvere} abandons element-centric containment. Instead, it defines a set entirely through its relational behavior using objects and functions (\textit{morphisms}). In this framework, we do not ask what an element \textit{is} internally; we ask how it behaves under functional composition. Lawvere’s Elementary Theory of the Category of Sets (\textbf{ETCS}) posits that the category of sets (\textbf{Set}) is a well-pointed, Cartesian Closed Category with a subobject classifier, a natural numbers object, and the axiom of choice.

To bridge this abstract mathematical formalism with systems-programming primitives, we translate the foundational axioms of ETCS into explicit C++23 type-predicates (\texttt{concept IsSet}) distributed across six localized compile-time partitions:

\subsection{1. Small Categories (\texttt{:small} partition)}
\begin{itemize}
    \item \textbf{Mathematical Axiom:} The existence of a collection of objects (sets) and arrows (functions/morphisms) that satisfy identity and associative composition laws.
    \item \textbf{C++ Reification:} Objects are bound to standard types satisfying \texttt{std::regular} to guarantee stable value identities. Morphisms are wrapped in custom templates satisfying \texttt{std::regular\_invocable<F, A>}.
    \item \textbf{Clarity Note for Engineers:} Instead of data-hiding containers, a ``set'' at this entry rung is simply a well-behaved C++ type acting as a domain, and a ``function'' is a side-effect-free transformation arrow mapping inputs to outputs. Identity arrows act as compile-time no-ops.
\end{itemize}

\subsection{2. Cartesian Categories (\texttt{:limit} partition)}
\begin{itemize}
    \item \textbf{Axioms of Terminal Object \& Finite Products:} 
    \begin{enumerate}
        \item \textit{Terminal Object:} There exists a unique set $1$ with exactly one element. Any set $A$ has exactly one function mapping $A \to 1$.
        \item \textit{Binary Products:} For any two sets $A$ and $B$, there exists a product set $A \times B$ equipped with projection arrows $\pi_1: A \times B \to A$ and $\pi_2: A \times B \to B$.
    \end{enumerate}
    \item \textbf{C++ Reification:} The terminal object is reified via \texttt{std::monostate}. Products are mapped to \texttt{std::pair} combined with compile-time structured bindings.
    \item \textbf{Clarity Note for Engineers:} This layer translates relational coordinate tables into type-safe tuples. The terminal object allows us to define ``elements'' of a set $A$ structurally as simple arrows from \texttt{std::monostate} to $A$.
\end{itemize}

\subsection{3. Cartesian Closed Categories (\texttt{:cartesian} partition)}
\begin{itemize}
    \item \textbf{Axiom of Exponentials:} For any sets $A$ and $B$, the collection of all functions from $A \to B$ forms a valid, manipulable set itself, denoted as $B^A$ (the exponential object).
    \item \textbf{C++ Reification:} Modeled via value-capturing \texttt{constexpr} lambda expressions and template instances of \texttt{std::function}.
    \item \textbf{Clarity Note for Engineers:} This licenses function \textit{Currying} entirely at compile time. It guarantees that functions are first-class citizens that can be passed, stored, and transformed by other functions without dropping out of the type-checking engine.
\end{itemize}

\subsection{4. Elementary Topoi (\texttt{:topoi} partition)}
\begin{itemize}
    \item \textbf{Axioms of Subobject Classifier \& Equalizers:}
    \begin{enumerate}
        \item \textit{Subobject Classifier:} There exists a truth-value set $\Omega$ and a unique arrow $\text{true}: 1 \to \Omega$. Every subset $S \subseteq A$ uniquely corresponds to a \textit{characteristic predicate} arrow $\chi: A \to \Omega$.
        \item \textit{Equalizers:} For any parallel functions $f, g: A \to B$, there exists a subset of $A$ where $f(x) = g(x)$.
    \end{enumerate}
    \item \textbf{C++ Reification:} Rather than forcing a hardcoded boolean infrastructure, $\Omega$ is generalized as a pluggable logical lattice type. It defaults to bit-perfect \texttt{bool} for classical two-valued logic, but can be swapped for three-valued structures like \texttt{struct K3} to represent undefined floating-point NaN states natively.
    \item \textbf{Clarity Note for Engineers:} An \textit{intensional} set is defined by a rule, not an array. Instead of searching a block of RAM to check if $x \in S$, the compiler executes the predicate function $\chi(x)$. If $\chi(x)$ evaluates to true, the element belongs to the subobject.
\end{itemize}

\subsection{5. Concrete Topoi (\texttt{:concrete} partition)}
\begin{itemize}
    \item \textbf{Axioms of the Powerset Monad \& Elementality:} Every set $A$ maps functorially to its powerset $\mathcal{P}(A)$ via a unit morphism $\eta_A: A \to \mathcal{P}(A)$.
    \item \textbf{C++ Reification:} Expressed via the \texttt{SetAsProduct} layout template, which combines a raw underlying C++ data type primitive with its corresponding logical evaluator lambda.
    \item \textbf{Clarity Note for Engineers:} This layer implements a faithful \textit{forgetful functor} $U$, which strips away structural categorical formatting to expose underlying raw C++ hardware types when evaluation is explicitly required. Constructing lazy evaluation pipelines (e.g., using \texttt{std::views::transform}) functions as monadic composition, chaining transformations without intermediate memory allocations.
\end{itemize}
Axiom 9: Power Set / Exponential Restriction (\(\beth _{1}\)). For any set X, there exists an object \(\mathcal{P}(X)\) classifying its subobjects. By Cantor’s theorem (Lawvere’s Fixed-Point Theorem), this explicitly bounds the cardinality of classical constructible sets in the base theory below the second transfinite cardinal. Russel's paradox is ruled out by construction. Library Reification: Restricts carrier predicates from modeling uncountably infinite domains beyond the continuum \(\beth _{1}\), avoiding type-system divergence and ensuring all logical intersections resolve during template translation.

\subsection{6. Set-Modeling Topoi (\texttt{:etcs} partition)}
\begin{itemize}
    \item \textbf{Axioms of Well-Pointedness, Infinity, and Choice:}
    \begin{enumerate}
        \item \textit{Well-Pointedness:} If two functions $f, g: A \to B$ are distinct, there must exist a single element $x: 1 \to A$ such that $f(x) \neq g(x)$.
        \item \textit{Infinity:} There exists a Natural Numbers Object (NNO) $\mathbb{N}$ that enables primitive recursion without structural loops.
        \item \textit{Choice:} Every surjective function has a right inverse section.
    \end{enumerate}
    \item \textbf{C++ Reification:} Enforced via the \texttt{IsSet} concept constraint. Well-pointedness enforces referential transparency and forbids hidden mutable states. Infinity is preserved by generating strict element-level partial orderings via \texttt{std::three\_way\_comparable} rather than managing physical loop limits.
    \item \textbf{Clarity Note for Engineers:} The NNO lets the library model \textit{transfinite spaces} (e.g., continuous mathematical boundaries or infinite sequence generators) lazily. The Axiom of Choice ensures that the template engine can constructively find or instantiate witness values for non-empty subset selections during compile-time verification passes.
\end{itemize}

